% Part: counterfactuals
% Chapter: minimal-change-semantics
% Section: introduction

\documentclass[../../../include/open-logic-section]{subfiles}

\begin{document}

\olfileid{cnt}{min}{int}

\olsection{పరిచయం}

``అగ్గిపుల్లను గీసి ఉంటే అది మండి
ఉండేది'' వంటి ప్రతివాస్తవ సోపాధికాలకు
స్టాల్నేకర్, లూయిస్ వివరణలను
ప్రతిపాదించారు. అటువంటి వాక్యాల సత్య
షరతులను ఎలా సరిగ్గా అర్థం చేసుకోవాలో
వారి ప్రతిపాదనలు వివరిస్తాయి. రెండు
ప్రతిపాదనల వెనుక భావం ఇదే:
ప్రతివాస్తవ సోపాధికం సత్యమా అని
నిర్ణయించాలంటే, పూర్వపక్షం సత్యమయ్యేలా
వాస్తవ లోకం ఉన్న తీరునుంచి కనిష్ఠంగా
మాత్రమే భిన్నమైన సాధ్యలోకాలను
పరిశీలించాలి. ఆ సాధ్యలోకాలలో
ఉత్తరపక్షం సత్యమైతే ప్రతివాస్తవ
సోపాధికం సత్యం. ఉదాహరణకు,
నా చేతిలో అగ్గిపుల్ల, అగ్గిపెట్టె
ఉన్నాయని అనుకుందాం. వాస్తవ లోకంలో
నేను వాటిని చూస్తూ, అగ్గిపుల్లను
గీస్తే ఏమవుతుందో ఆలోచిస్తున్నాను.
నేను అగ్గిపుల్లను గీసే లోకంలో,
వాస్తవ లోకానికి కనిష్ఠ మార్పు
ఏమిటంటే నేను అలా చేయాలని
నిర్ణయించి దాన్ని గీయడమే.
మరేదీ మారదు కాబట్టి అది కనిష్ఠం:
నేను గాలిలోకి ఎగరను; అగ్గిపుల్ల
గీయడంవల్ల నా జుట్టుకు కూడా
మంట అంటుకోదు; నా వేళ్లలోని
శక్తి అకస్మాత్తుగా పోదు; అదే
సమయంలో నీటి తుపాకీ దాడిలో
నేను తడిచిపోను; వగైరా.
ఆ ప్రత్యామ్నాయ సాధ్యతలో
అగ్గిపుల్ల మండుతుంది. కనుక
అగ్గిపుల్లను గీసి ఉంటే అది
మండి ఉండేదనేది సత్యం.

ఈ సహజ అవగాహనకు ప్రతివాస్తవ
తర్కాల నియత అర్థవిచారాన్ని
జోడించవచ్చు. లూయిస్ ప్రతివాస్తవానికి
``$\boxright$'' సంకేతాన్ని,
స్టాల్నేకర్~``$>$'' సంకేతాన్ని
ప్రవేశపెట్టారు. మనం~$\cif$ను
వాడి, ప్రతిజ్ఞా తర్కానికి ద్విస్థానిక
సంయోజకంగా చేరుస్తాం. అప్పుడు
$!A \lif !B$ రూపపు
\tetoken{సూత్రాలతో}{!!{formula}s}
పాటు~$!A \cif !B$ రూపపు
\tetoken{సూత్రాలూ}{!!{formula}s}
ఉంటాయి. మోడల్ తర్క సంబంధ
అర్థవిచారం మాదిరిగానే, ఈ నియత
అర్థవిచారం లోకాల వద్ద
\tetoken{సూత్రాలను}{!!{formula}s}
మూల్యాంకనం చేసే నమూనాలపై ఆధారపడుతుంది.
$\mSat{M}{!A \cif !B}[w]$ను
నిర్వచించే సంతృప్తి షరతు, ఇతర
కొన్ని లోకాల~$w'$ వద్ద
$\mSat{M}{!A}[w']$,
$\mSat{M}{!B}[w']$లపై ఆధారపడుతుంది.
ఏ $w'$లు? సహజంగా,~$w$కు
అత్యంత సమీపంగా ఉండి
~$\mSat{M}{!A}[w']$ అయ్యేవి.
అందుకోసం ``సమీపత్వ'' సంబంధాన్నీ
నమూనాలో చేర్చాలి.

ప్రతివాస్తవ పరిస్థితులను చిత్రంగా
చూపే బోధప్రదమైన పద్ధతిని లూయిస్
ప్రవేశపెట్టారు. ప్రతి సాధ్యలోకం,
ఇతర లోకాలను కలిగిన ఒకదానిలో
మరొకటి ఉన్న గోళాల సమితికి
కేంద్రంగా ఉంటుంది---ఈ గోళాలను
ఏకకేంద్ర వృత్తాలుగా గీస్తాం.
రెండు గోళాల మధ్యనున్న లోకాలు
కేంద్ర లోకానికి ఒకదానితో మరొకటి
సమాన సమీపంలో ఉంటాయి; అంతర్గత
గోళంలోని లోకాలు మరింత దగ్గరగా,
బాహ్య గోళంలోని లోకాలు మరింత
దూరంగా ఉంటాయి.
\begin{center}
\begin{tikzpicture}[scale=.7]\small
  \spheresystem{5}
  \draw (0,0) node {$w$};
  \propositionintersect{0}{4}{40}{3.5}
  {\tikzset{proposition/.append={smooth,tension=1.4}}
  \proposition[shift={(1.6,-1)}]{90}{1}{30}{4}}
  \path (1.5,3) node[below] {$\formula{B}$};
  \path (3,0) node {$\formula{A}$};
\end{tikzpicture}
\end{center}
అత్యంత సమీప $!A$-లోకాలు
అంటే కేంద్ర లోకం~$w$ చుట్టూ
ఉన్న అత్యంత చిన్న గోళంలో
(బూడిద రంగు ప్రాంతంలో)
ఉండి,~$!A$ సంతృప్తమయ్యే
లోకాలు~$w'$. సహజ అవగాహన
ప్రకారం, అత్యంత సమీప
$!A$-లోకాలన్నిటిలో $!B$
సత్యమైతే~$w$ వద్ద
$!A \cif !B$ సంతృప్తమవుతుంది.

\end{document}
